\documentclass[11pt,reqno]{amsart}
\usepackage[T1]{fontenc}
\usepackage{lmodern,amsmath,amssymb,mathtools,microtype,booktabs,array,longtable}
\usepackage[margin=1in]{geometry}
\usepackage[hidelinks]{hyperref}
\ifdefined\pdfinfoomitdate\pdfinfoomitdate=1\fi
\ifdefined\pdftrailerid\pdftrailerid{}\fi
\ifdefined\pdfsuppressptexinfo\pdfsuppressptexinfo=15\fi
\numberwithin{equation}{section}
\newtheorem{theorem}{Theorem}[section]
\newtheorem{proposition}[theorem]{Proposition}
\newtheorem{lemma}[theorem]{Lemma}
\newtheorem{corollary}[theorem]{Corollary}
\theoremstyle{definition}\newtheorem{definition}[theorem]{Definition}
\newtheorem{example}[theorem]{Example}
\theoremstyle{remark}\newtheorem{remark}[theorem]{Remark}
\newcommand{\E}{\mathbb E}\newcommand{\Prb}{\mathbb P}
\newcommand{\R}{\mathbb R}\newcommand{\one}{\mathbf1}
\newcommand{\Q}{\mathcal Q}\newcommand{\TV}{\mathrm{TV}}
\newcommand{\cp}{\mathrm{cp}}\newcommand{\ev}{\mathrm{ev}}\newcommand{\st}{\mathrm{st}}
\DeclareMathOperator{\tr}{tr}\DeclareMathOperator{\dist}{dist}
\DeclareMathOperator{\law}{Law}\DeclareMathOperator{\rank}{rank}
\DeclareMathOperator{\supp}{supp}
\emergencystretch=2em

\title[Technical supplement]{General Theta Foundations I:\\Technical Supplement to Acquired Geometry and Causal Resource Transfer}
\author{Qian Qi}\date{October 2026}
\hypersetup{pdftitle={General Theta Foundations I: Current Technical Supplement},pdfauthor={Qian Qi}}
\begin{document}
\begin{abstract}
This supplement supplies the preparation-face, legal-domain, deadlock and post-report compilation arguments used by the complete current manuscript. It distinguishes termination supports from report-resolved task spans and survival probabilities from expected cost tails. The companion on absolute calibration, stopping resources and nonintervening sensing is separately retained in full. The statements here are mathematical proofs with explicit scopes, not formal verification or editorial certification.
\end{abstract}
\maketitle

\section{Preparation faces and reachable supports}
Let $\mathcal A=\bigoplus_{s=1}^{\ell}M_d(\mathbb C)$ be the live controller--physical algebra. A fixed controller has a completely positive trace-nonincreasing continuation map $\mathcal T$, a cost observable $g\succeq c_*I$ and terminal score observables $f_j$ satisfying
\[
 \alpha_j(I-\mathcal T^*I)\preccurlyeq f_j\preccurlyeq
 \beta_j(I-\mathcal T^*I).
\]
The initial live label is $s_0$. A preparation face is specified by a nonzero physical projection $Q$, and all initial density matrices supported on $Q$ are admitted. Let $\iota$ insert a matrix into block $s_0$, and put
\[
 X_0=\iota(Q),\quad D=\ell d,\quad R=\sum_{n=0}^{D-1}\mathcal T^nX_0,
 \quad P=\supp R.
\]
\begin{theorem}[Exact face-restricted potential certificate]\label{thm:face}
The projection $P$ is the least invariant block support containing $X_0$. A fixed controller is proper at expected cost at most $N$ for every admitted preparation, with prescribed terminal operators $A_j$ on $Q$, if and only if there exist an invariant block projection $\widetilde P$, with $\widetilde P X_0=X_0$, a positive $H=\widetilde P H\widetilde P$, and Hermitian $K_j=\widetilde P K_j\widetilde P$ such that
\begin{align*}
 \widetilde P(H-\mathcal T^*H-g)\widetilde P&\succeq0,
 &QH_{s_0}Q&\preccurlyeq NQ,\\
 \widetilde P(K_j-\mathcal T^*K_j-f_j)\widetilde P&=0,
 &Q(K_j)_{s_0}Q&=A_j,\\
 \alpha_j\widetilde P&\preccurlyeq K_j\preccurlyeq\beta_j\widetilde P.
\end{align*}
One may choose $\widetilde P=P$. On that least support, finite-cost proper operation from the preparation face is equivalent to $\operatorname{spr}(\mathcal T_P)<1$, and the exact minimal potentials are the supported Neumann resolvents. No transience is required on an unreachable block.
\end{theorem}
\begin{proof}
The supports of $R_n=\sum_{k=0}^n\mathcal T^kX_0$ increase. When two successive supports agree, $R_n$ dominates a positive multiple of its support projection. Positivity then implies that the support is invariant under $\mathcal T$. Until that event the rank increases, and its ambient bound is $D$; hence stabilization occurs by $D-1$. Every invariant support containing $X_0$ contains every iterate, proving leastness.

For sufficiency, iterate the cost superpotential inequality on $\widetilde P$ and pair with an initial density matrix supported on $Q$ in block $s_0$. Dropping the positive remainder bounds every partial expected cost by $N$. Monotone convergence bounds total cost. The positive floor bounds $c_*\sum_{k\ge0}\Pr(\tau>k)$ by $N$, proving properness. Iterate the score equation. Its terminal remainder is bounded in absolute trace by $\max(|\alpha_j|,|\beta_j|)\Pr(\tau>n)$, which vanishes. Its partial sum contains the actual disjoint terminal-score contributions. Thus the initial compressed block is exactly $A_j$.

Conversely the cost bound on every initial density in the face gives
$\sum_{k\ge0}\langle g,\mathcal T^kX_0\rangle\le N\rank Q$.
The same series starting at $\mathcal T^nX_0$ is a nonnegative tail of that series. Summing for $0\le n<D$ yields a finite bound $DN\rank Q$ on the series starting at $R$. Since $R$ is strictly positive on its support, every positive supported matrix is bounded by a finite multiple of $R$. Thus the increasing partial cost potentials on that support have bounded norm and converge. Their limit $H^{\min}$ solves $H^{\min}-\mathcal T_P^*H^{\min}=PgP$ and has the required compressed initial budget.

The floor implies norm convergence of $\sum_k(\mathcal T_P^*)^kP$. Positivity bounds the iterates of any supported Hermitian matrix by its norm times the iterates of $P$, so every such iterate tends to zero. The same holds for all complex matrices by decomposing into Hermitian parts. In finite dimension this forces every eigenvalue of $\mathcal T_P^*$ to have modulus less than one, proving transience and the Neumann formula. The bounded-score series also converges. Its normalization bounds follow from telescoping the terminal effect $P-\mathcal T_P^*P$, whose series sums to $P$. The initial face identity follows from the assumed terminal scores. Finally iteration shows that every positive superpotential dominates every partial exact cost potential, hence their limit.
\end{proof}

An arbitrary preparation set is not automatically its whole support face. For equality of task expectations, its linear span, not its support face, controls what is required. The main manuscript's report-resolved continuation theorem treats that distinction. Extending a preparation family to its full face strengthens the quantifier and must be declared explicitly.

\section{Legal physical domains and deadlocks}
A physical command can have a restricted domain. A time-homogeneous face-valued domain is represented by projections $D_a$ on the physical space: command $a$ is permitted only on states supported on $D_a$. The controller is not given that support or a free physical-port tag. Let $P_s$ be the controller's least reached physical support at label $s$, computed from the allowed preparations and the actual continuing instrument.
\begin{proposition}[Fixed-row legality on reached supports]\label{prop:legal}
A command row $p_{s,a}$ is legal at every reached history if and only if
\[
 p_{s,a}>0\quad\Longrightarrow\quad P_s\preccurlyeq D_a
\]
for every label and command. Equivalently $p_{s,a}(I-D_a)P_s=0$. For a finite algebraic model family with face domains, these are polynomial conditions on the shared row probabilities and modelwise support projections. They may be added to the finite-memory potential system without giving the controller an extra domain input.
\end{proposition}
\begin{proof}
If the support inclusion holds, every reached positive state at label $s$ is supported on $D_a$, so the command is legal. Conversely, legality at every reached positive state annihilates all of them by $I-D_a$ whenever $p_{s,a}>0$. Their positive sum over a spanning finite set of reached supports has support $P_s$; hence $I-D_a$ also annihilates $P_s$. This gives the inclusion. The scalar multiplication formula handles zero row probabilities without introducing a division or a nonclosed strict-positivity condition. It is polynomial in real and imaginary matrix coordinates. Apply it separately in each model, retaining the same row probabilities.
\end{proof}

This proposition concerns declared support-face domains. A legal domain specified by a different predicate on states need not admit this polynomial encoding. If a protocol port is already a deterministic function of the visible history, a full-history policy can recover it, but a finite-label controller must store enough information to choose a legal row without a free port input.

For the full-history finite-obstruction hierarchy, the main manuscript assumes a nonempty legal command set at every open decision node. Deadlocks can instead be admitted explicitly as forbidden nonproductive leaves. A node with no legal decision must carry zero reached mass. For a finite set of models and finite horizon, impose the corresponding positive-operator equality for each such leaf on the admitted preparation face. In the common realization plan it is the equality
\[
 Q_\theta\,\gamma(h)\mathcal I_{\theta,h}^*(I)\,Q_\theta=0.
\]
Here $\gamma(h)$ excludes physical report probabilities, as in the main manuscript.
\begin{proposition}[Deadlock-aware finite hierarchy]\label{prop:deadlock}
With finitely branching legal histories, a common initial preparation face in each finite-dimensional model, a positive call-cost floor and bounded terminal scores, the finite-intersection proof of the causal synthesis hierarchy remains valid after adding the deadlock zero-mass constraints. Constraints at depth $T$ include every deadlocked leaf reached by that depth. An open nondeadlocked leaf still owes at least one further paid call. A deadlocked leaf cannot be used as a fictitious terminal output or a free fallback.
\end{proposition}
\begin{proof}
Give an empty decision node a singleton dummy coordinate solely to make the product of nodewise controller coordinate spaces compact. This coordinate causes no physical action and is constrained to have zero reached mass in each admitted model. It is not a legal executed command. All finite-prefix mass, cost and score expressions remain continuous in the product coordinates. Deadlock constraints are closed and remain true under extension to later horizons. Restricting a later feasible prefix therefore preserves them and the usual cost reserve and score intervals. Conversely, every proper legal policy satisfies them because it never reaches a deadlock with positive probability.

For any finite collection of model, horizon and deadlock constraints, choose its largest horizon and union of model indices. Feasibility at that level gives one common prefix satisfying the collection. Compactness gives an infinite common coordinate assignment satisfying all levels. No dummy decision is executed with positive probability. The same positive-cost survival bound tends to zero, so the resulting process commits almost surely without reaching a deadlock, and monotone convergence supplies its budget. Bounded terminal-score tails recover its targets. This repeats the hierarchy argument with the missing legal-boundary event explicitly accounted for.
\end{proof}

\section{Prospective audits and post-report return compilation}
A routine's logical return and the experiment's irrevocable halt are different events. Let its final productive call produce a visible report $y$. At this report the routine selects a return mark $x$ and a finite exported decoder choice, using current coins. Its prospective audit instrument defines the score that would be obtained if the outer caller accepted. Those scores do not require executing the audit on a rejected branch.
\begin{proposition}[Compilation without a free audit or hidden mode]\label{prop:compile}
Suppose the caller's acceptance probability and decoder are functions only of the current finite return mark and the existing finite decoder dictionary. Suppose the physical rejection operation is legal and either executes within the current productive call or has an explicitly counted finite-state return subroutine. Then the routine can be compiled into an event-stopped observer by replacing each logical return transition by its acceptance and rejection branches. On acceptance it emits the outer halt event and retains the decoder label; on rejection it executes the declared return and reenters the existing start label. Prospective audit statistics on accepted branches, complete physical paths and charged costs agree with this construction. The state-stopped implementation uses disjoint live and final-output types instead.
\end{proposition}
\begin{proof}
At the current productive report, first sample exactly the routine's original return choice. Sample the caller's independent acceptance coin next. On acceptance, emit the outer commitment event, store its dictionary value and only then perform the declared scored audit. On rejection do not perform that audit; follow the declared rejection instrument and reenter the start. This reproduces each branch's raw physical instrument and its unnormalized accepted audit measure. Iteration gives equality of every finite stopped prefix and every raw path cost. Properness is supplied by the positive conditional-hazard argument, not by the compilation alone.

All information used to decide the branch is available during the current call. No value of the future audit is read. If the return needs another physical call, its phase persists and is stored in the added return-subroutine labels; its cost is included. Otherwise the existing start label suffices for continuation, and the event-stopped live/output packing introduces no continuing mode bit. A state-stopped implementation sends accepted branches to disjoint halt-typed labels. This proves the assertions while keeping the different observer models distinct.
\end{proof}

A synthesized terminal score list alone is not enough to prove this proposition's hypotheses. In particular, a score obtained only by actually measuring and revealing the audit cannot be treated as an unobserved prospective score on rejection. Likewise a kernel identity at the first routine root cannot substitute for its report-resolved continuation certificate. The current main manuscript supplies that additional finite test for time-homogeneous finite-return routines.

\section{Survival mass is not expected cost tail}
\begin{proposition}[The exact limitation of a mean-cost bound]\label{prop:tail}
For a positive-integer stopping time with $\E\tau\le b$, one has
\[
 \Pr(\tau>T)\le\min\{1,(b-1)/T\}\quad(T\ge1).
\]
This does not imply a uniform vanishing bound on $\E[(\tau-T)_+]$ over the same budget class. In particular the main manuscript's score-tail convergence, which uses bounded scores, cannot be repurposed as an expected cost-tail convergence statement.
\end{proposition}
\begin{proof}
Pointwise $\tau\ge1+T\one_{\{\tau>T\}}$, which gives the probability bound. Fix $b>1$. For an integer $K\ge b$, set $\tau=K$ with probability $(b-1)/(K-1)$ and $\tau=1$ otherwise. This is a finite-branching, unit-call-cost terminating experiment and $\E\tau=b$. At a fixed cutoff $T<K$,
\[
 \E[(\tau-T)_+]=(b-1)\frac{K-T}{K-1}\longrightarrow b-1
 \quad(K\longrightarrow\infty).
\]
Thus its supremum over the budget class never tends to zero with $T$. The survival probability at that same cutoff can tend to zero. Bounded terminal scores are controlled by the latter probability; unbounded remaining call cost is not.
\end{proof}

The finite-prefix fallback theorem uses a conditional cost bound on an actual fallback, not an unproved uniform tail-cost estimate. For a fixed finite-state transient instrument its own supported resolvent gives a finite tail-cost bound, but its constants need not be uniform over a family approaching the boundary of transience.

\section{Correction boundary and preserved companion}
The single-atom proposition inherited by the current manuscript is corrected only in its necessary-threshold quantifier. The statement that its specified atom yields zero loss when $b\lambda\ge1$ remains valid, as do its lower and upper comparisons with the remainder profile. Necessity based on that atom alone requires positive residual variance almost everywhere on the remainder, or another condition excluding all other finite-dictionary zero-loss mass. The current main theorem gives the complete criterion through the largest atomic masses of the zero-residual predictive measure, and an explicit two-atom counterexample shows why the qualification matters.

The byte-preserved R58 native companion has its own older statements and provenance. Retention is not a claim that its historical wording supersedes a correction in the current manuscript. The preservation map identifies the current controlling result for every corrected claim. The original acquired-law geometry, absolute-calibration proofs, nonintervening Gaussian sensing and exact stopping constructions remain available in that companion, and the independent ordered-measurement, reset-communication and rank-constrained discrimination developments remain at their pinned source identities.

\section{Verification status}
The accompanying scripts compare finite algebraic examples, solve small finite reset-control instances independently of the manuscript's closed formulas, check report-resolved span counterexamples, test boundary resource values, and compile the complete native manuscript, this supplement and the retained companion. Exact rational checks and floating-point residual checks are separately counted. Neither kind proves an unbounded theorem by sampling. Native LaTeX compilation verifies source completeness and references, not mathematical correctness. No formal proof assistant, newly commissioned independent human referee or editorial acceptance is claimed.
\end{document}
